\section{视觉的起源：从寒武纪到人类文明}

\subsection{5.4 亿年前的寒武纪大爆发}

视觉的历史并非始于人类文明，而是可以追溯到 \textbf{5.4 亿年前}的寒武纪大爆发（Cambrian Explosion）。化石研究显示，在大约 1000 万年的进化时间窗口内（对进化而言极为短暂），地球上的动物物种突然大量涌现。

\begin{figure}[H]
\centering
\includegraphics[width=0.95\textwidth]{slides-images/slide-012.jpg}
\caption{寒武纪大爆发：5.4亿年前三叶虫等动物获得了最早的感光细胞，视觉的出现可能是物种大爆发的关键驱动力\protect\footnotemark}
\end{figure}
\footnotetext{Slides 第13页（Part 1）。}

\begin{knowledgebox}{视觉与智能的关系}
寒武纪大爆发最有说服力的理论之一是\textbf{视觉的出现}。最早的动物——三叶虫（trilobite）获得了感光细胞（photosensitive cells），这并非精密的视网膜和透镜系统，而仅仅是一个能收集光线的小孔（pinhole）。然而，这个简单的能力彻底改变了生命的形态：
\begin{itemize}
    \item 没有感觉的生命只是被动的新陈代谢
    \item 有了视觉，动物成为环境的积极参与者——觅食、躲避天敌、竞争资源
    \item 进化压力驱动了神经系统和智能的发展
\end{itemize}
Fei-Fei Li 总结道：``5.4 亿年的视觉进化史，就是智能的进化史。''
\end{knowledgebox}

人类是高度视觉化的动物：超过一半的皮层神经元参与视觉处理。我们拥有非常复杂的视觉系统，这也是研究视觉智能如此令人兴奋的原因。

\subsection{从暗箱到数字相机}

人类文明很早就开始尝试制造``看''的机器。从古希腊和古代中国关于针孔成像的思考，到达芬奇对暗箱（Camera Obscura）的研究，再到现代数字相机的爆发——硬件层面的``眼睛''已经无处不在。

\begin{figure}[H]
\centering
\includegraphics[width=0.95\textwidth]{slides-images/slide-015.jpg}
\caption{暗箱（Camera Obscura）的历史：从 Gemma Frisius（1545）到达芬奇（16世纪）再到18世纪百科全书中的设计\protect\footnotemark}
\end{figure}
\footnotetext{Slides 第16页（Part 1）。}

\begin{warningbox}{相机不等于视觉}
相机只是采集光信号的装置，就像眼睛只是光学器官一样。真正的挑战在于\textbf{理解}视觉信息——这才是计算机视觉要解决的核心问题。正如 Fei-Fei Li 所说：``Cameras are not enough for seeing, just like eyes are not enough for seeing.''
\end{warningbox}

\subsection{视觉的根本难题：从 2D 到 3D 的逆问题}

视觉的一个根本难题在于：真实世界是三维的，但视觉感知到的（无论是视网膜还是相机传感器）是二维的投影。从二维图像恢复三维世界的完整信息，在数学上是一个\textbf{不适定问题}（ill-posed problem）。

自然界的解决方案之一是发展出多只眼睛（大多数动物是两只），通过三角测量来估计深度。但仅靠两只眼睛远远不够，还需要理解对应关系（correspondences）等复杂计算。

\begin{knowledgebox}{语言与视觉的哲学差异}
Fei-Fei Li 提出了一个深刻的观察：\textbf{语言不存在于自然界中}——它是人类大脑生成的、一维的、序列化的产物。这也是为什么大语言模型（LLM）如此成功——语言的结构天然适合序列建模。而\textbf{视觉则植根于物理世界}，受物理定律、材料属性等约束，这使得视觉问题具有本质上不同的挑战。
\end{knowledgebox}

\subsection{本章小结}

视觉是生命进化的核心驱动力，从5.4亿年前的寒武纪大爆发到现代数字相机的普及，人类对``看''的追求从未停止。但真正的挑战不在于采集图像，而在于理解图像——从二维投影中恢复三维世界的语义信息。这是一个数学上不适定的、计算上极具挑战性的问题。

%% ========== 计算机视觉简史 ==========

